\documentclass[../../main.tex]{subfiles}

\begin{document}

\chapter[Bottom-Up Parsing]{Bottom-Up Parsing: Shift-Reduce\\ dan Operator Precedence}
\label{chap:bottom-up}

\begin{outcomebox}
\textbf{Sub-CPMK pendukung:} \subcpmkref{4.2}\\[3pt]
Setelah menyelesaikan bab ini, mahasiswa mampu:
\begin{enumerate}[label=\arabic*.,leftmargin=*]
  \item menjelaskan model shift-reduce dan peran stack;
  \item mengidentifikasi handle dan viable prefix;
  \item menjelaskan konflik shift-reduce dan reduce-reduce;
  \item membangun tabel operator precedence sebagai fondasi;
  \item menyiapkan data untuk tabel ACTION/GOTO SLR/LALR yang
        dibangun pada Bab XII.
\end{enumerate}
\end{outcomebox}

\section*{Materi Pokok}
\begin{itemize}[leftmargin=*]
  \item Model shift-reduce, stack, dan handle.
  \item Viable prefix dan konflik.
  \item Operator precedence grammar dan tabelnya.
  \item Algoritma operator precedence dan deteksi kesalahan.
  \item Jembatan dari operator precedence ke tabel ACTION/GOTO
        SLR/LALR.
\end{itemize}

\subfile{section-01}
\subfile{section-02}
\subfile{section-03}
\subfile{section-04}
\subfile{section-05}

\section*{Contoh}
\addcontentsline{toc}{section}{Contoh}

\begin{contoh}[Menemukan handle]
Pada bentuk kalimat \(\code{id} + \code{id} * \code{id}\), handle
pertama adalah \code{id} paling kiri. Setelah direduksi menjadi
\(F\), \(T\), dan \(E\), handle berikutnya adalah \code{id} kedua.
Pemilihan handle yang salah, misalnya mereduksi seluruh
\(\code{id} + \code{id}\) terlebih dahulu, menghasilkan pohon yang
tidak sesuai precedence.
\end{contoh}

\begin{contoh}[Relasi precedence]
Untuk \(\code{id} + \code{id} * \code{id}\):
\[
\code{id} \cdot> +, \quad
+ <\cdot \code{id}, \quad
+ <\cdot *, \quad
* \cdot> \code{id}.
\]
Relasi ini memaksa perkalian direduksi sebelum penjumlahan.
\end{contoh}

\begin{contoh}[Unary minus]
Operator unary minus memiliki precedence lebih tinggi daripada
binary minus. Pada ekspresi \(-a * b\), unary minus diterapkan pada
\(a\) terlebih dahulu, lalu hasilnya dikalikan \(b\). Grammar harus
membedakan aturan unary dan binary, atau parser generator harus
mendeklarasikan precedence unary secara eksplisit.
\end{contoh}

\begin{contoh}[Kesalahan pada operator precedence]
Pada masukan \(\code{id} + + \code{id}\), setelah shift \(+\) kedua,
tabel tidak memiliki relasi antara \(+\) dan \(+\). Parser
melaporkan kesalahan pada token kedua. Pemulihan dapat dilakukan
dengan melewati token sampai operator atau operand berikutnya.
\end{contoh}

\begin{contoh}[Membaca keluaran kalkulator]
Program \file{shift\_reduce.c} mencetak urutan shift dan reduce.
Untuk masukan \code{2+3*4}, program mencetak shift angka 2, shift
operator +, shift angka 3, shift operator *, shift angka 4, lalu
reduce 3*4, reduce 2+12, dan hasil 14. Urutan ini menunjukkan bahwa
perkalian direduksi lebih dahulu.
\end{contoh}

\section*{Aktivitas Pembelajaran dan Praktikum}
\addcontentsline{toc}{section}{Aktivitas Pembelajaran dan Praktikum}

\begin{table}[H]
\centering
\caption{Rencana aktivitas Bab XI}
\label{tab:ch11-aktivitas}
\begin{tabularx}{\textwidth}{@{}llX@{}}
\toprule
\textbf{Sesi} & \textbf{Bentuk} & \textbf{Kegiatan} \\
\midrule
1 & Ceramah & Model shift-reduce dan handle \\
2 & Latihan & Menelusuri stack untuk ekspresi \\
3 & Lokakarya & Membangun tabel operator precedence \\
4 & Praktikum & Kalkulator operator precedence \\
\bottomrule
\end{tabularx}
\end{table}

\subsection*{Lokakarya: Penelusuran dan Kalkulator}
\textbf{Tujuan.} Memahami shift-reduce melalui penelusuran manual
dan implementasi kalkulator sederhana.

\textbf{Langkah kerja.}
\begin{enumerate}[label=\arabic*.,leftmargin=*]
  \item Telusuri shift-reduce untuk \(\code{id} * \code{id} +
        \code{id}\).
  \item Tandai handle pada setiap langkah reduce.
  \item Bangun tabel operator precedence untuk operator
        \(+\), \(-\), \(*\), \(/\), dan tanda kurung.
  \item Modifikasi \file{shift\_reduce.c} agar mendukung operator
        \(-\), \(/\), dan tanda kurung.
  \item Uji dengan ekspresi yang mengandung precedence dan
        associativity berbeda.
\end{enumerate}

\textbf{Definisi selesai.}
\begin{checklist}
  \item Penelusuran manual benar.
  \item Tabel precedence lengkap.
  \item Kalkulator mendukung empat operator dan kurung.
  \item Minimal 10 ekspresi uji.
  \item Data penelusuran dan tabel precedence siap dipakai pada
        Bab XII.
\end{checklist}

Bab ini tidak menambahkan milestone proyek baru. Hasil lokakarya
menjadi data pendukung Praktikum 3 (12\%) yang dilaporkan pada
Bab XIII. Tabel ACTION/GOTO SLR/LALR dibangun pada Bab XII.

\section*{Latihan dan Refleksi}
\addcontentsline{toc}{section}{Latihan dan Refleksi}

\subsection*{Latihan}
\begin{enumerate}[label=\arabic*.,leftmargin=*]
  \item Jelaskan perbedaan shift dan reduce.
  \item Definisikan handle dan berikan contoh.
  \item Telusuri shift-reduce untuk \(\code{id} * \code{id} +
        \code{id}\).
  \item Jelaskan viable prefix.
  \item Berikan contoh konflik shift-reduce.
  \item Berikan contoh konflik reduce-reduce.
  \item Bangun tabel operator precedence untuk \(+\), \(*\), dan
        \(\code{id}\).
  \item Bandingkan top-down dan bottom-up dari sisi grammar yang
        didukung.
\end{enumerate}

\subsection*{Refleksi}
\begin{reflectionbox}
\begin{enumerate}[label=\arabic*.,leftmargin=*]
  \item Apakah Anda dapat menjelaskan mengapa bottom-up menangani
        left recursion?
  \item Bagian mana dari penelusuran shift-reduce yang paling sulit?
  \item Mengapa konflik perlu diselesaikan secara eksplisit?
\end{enumerate}
\end{reflectionbox}

\section*{Asesmen}
\addcontentsline{toc}{section}{Asesmen}

Asesmen Bab XI menjadi bagian UTS (20\%) dan Praktikum 3 (7\%).
Fokusnya adalah penelusuran shift-reduce dan operator precedence.

\begin{table}[H]
\centering
\caption{Instrumen asesmen Bab XI}
\label{tab:ch11-asesmen}
\begin{tabularx}{\textwidth}{@{}lXl@{}}
\toprule
\textbf{Instrumen} & \textbf{Indikator} & \textbf{Bobot} \\
\midrule
UTS & Penelusuran shift-reduce dan konflik & 20\% \\
Lokakarya & Tabel precedence dan kalkulator & bagian P3 \\
Kuis & Handle dan viable prefix & formatif \\
\bottomrule
\end{tabularx}
\end{table}

Kriteria kelulusan: mahasiswa dapat menelusuri shift-reduce,
menjelaskan handle dan konflik, serta membangun tabel operator
precedence sederhana.

\section*{Checklist Pencapaian Kompetensi}
\addcontentsline{toc}{section}{Checklist Pencapaian Kompetensi}

\begin{checklist}
  \item Saya dapat menjelaskan model shift-reduce.
  \item Saya dapat mengidentifikasi handle.
  \item Saya dapat menjelaskan viable prefix.
  \item Saya dapat menjelaskan konflik shift-reduce dan
        reduce-reduce.
  \item Saya dapat membangun tabel operator precedence.
  \item Saya dapat menelusuri parser operator precedence.
  \item Saya dapat menjelaskan keterbatasan operator precedence dan
        kapan parser LR diperlukan.
  \item Saya dapat menyiapkan data untuk tabel ACTION/GOTO SLR/LALR
        pada Bab XII.
  \item Saya dapat membandingkan top-down dan bottom-up.
  \item Saya sudah menyelesaikan lokakarya kalkulator.
\end{checklist}

\section*{Rangkuman}
\addcontentsline{toc}{section}{Rangkuman}

Pada bab ini telah dibahas:
\begin{itemize}[leftmargin=*]
  \item model shift-reduce, stack, dan handle;
  \item viable prefix dan konflik shift-reduce/reduce-reduce;
  \item operator precedence grammar dan tabelnya;
  \item algoritma operator precedence dan deteksi kesalahan;
  \item data yang disiapkan untuk tabel ACTION/GOTO SLR/LALR;
  \item perbandingan parsing top-down dan bottom-up.
\end{itemize}

Bab XII akan memperluas bottom-up parsing menjadi LR parsing:
item LR(0), SLR(1), canonical LR(1), dan LALR(1). LR menyediakan
tabel sistematis dan menangani grammar yang lebih luas daripada
operator precedence.

\section*{Evaluasi}
\addcontentsline{toc}{section}{Evaluasi}

\begin{enumerate}[label=\arabic*.,leftmargin=*]
  \item Jelaskan perbedaan shift dan reduce.
  \item Definisikan handle dan viable prefix.
  \item Telusuri shift-reduce untuk \(\code{id} + \code{id} *
        \code{id}\).
  \item Jelaskan konflik shift-reduce dengan contoh.
  \item Jelaskan konflik reduce-reduce dengan contoh.
  \item Bangun tabel operator precedence untuk \(+\), \(*\),
        \(\code{id}\), dan \(\$\).
  \item Mengapa operator precedence tidak cocok untuk seluruh bahasa?
  \item Bandingkan kelebihan top-down dan bottom-up.
  \item Jelaskan data apa yang perlu ditambahkan agar tabel operator
        precedence dapat diperluas menjadi tabel ACTION/GOTO SLR/LALR.
\end{enumerate}

\subsection*{Kunci dan Rubrik Singkat}
\begin{itemize}[leftmargin=*]
  \item Soal 1--2: ketepatan definisi (70\%), contoh (30\%).
  \item Soal 3--5: kebenaran penelusuran dan konflik (70\%),
        penjelasan (30\%).
  \item Soal 6--9: ketepatan tabel dan analisis (60\%),
        argumentasi (40\%).
\end{itemize}


\ifSubfilesClassLoaded{%
  \bibliographystyle{apalike}%
  \bibliography{../../references}%
}{}

\end{document}
